% 图 4.1 TVM 编译流水线（重绘加强版）
% 基于 CICC0903540 技术文档 4.1 - TVM 优化编译器框架
\documentclass[tikz,border=10pt]{standalone}

% 强制全局样式与配色
\usepackage{amssymb, amsmath}
\usepackage{fontspec}
\usepackage{xeCJK}
\setCJKmainfont{SimHei}
\usepackage{tikz}
\usetikzlibrary{
  arrows.meta, positioning, shapes.geometric, calc,
  backgrounds, fit, decorations.pathreplacing, patterns, shapes.symbols
}

% 低饱和度马卡龙配色
\definecolor{mBlue}{HTML}{AEC6CF}
\definecolor{mPink}{HTML}{FFD1DC}
\definecolor{mGreen}{HTML}{B1D8B7}
\definecolor{mPurple}{HTML}{B39EB5}
\definecolor{mYellow}{HTML}{FDFD96}
\definecolor{mGray}{HTML}{E5E4E2}
\definecolor{mDark}{HTML}{4A4A4A}

% 全局 TikZ 样式
\tikzset{
  >=stealth,
  thick,
  draw=mDark,
  stage/.style={
    rectangle, rounded corners=3pt, draw=mDark, align=center,
    font=\sffamily\small, minimum height=0.9cm, inner sep=5pt,
    minimum width=2.2cm},
  micro/.style={
    rectangle, rounded corners=2pt, draw=mDark!70, align=center,
    font=\sffamily\tiny, minimum height=0.4cm, inner sep=3pt},
  ts/.style={font=\tiny\sffamily, color=mDark!80},
  arr/.style={-{Stealth[length=4pt]}, thick, draw=mDark},
  legendbox/.style={
    rectangle, draw=mDark!50, rounded corners=3pt,
    fill=mGray!15, inner sep=5pt, font=\tiny\sffamily},
}

\begin{document}
\begin{tikzpicture}[
  node distance=0.6cm and 0.8cm,
]
  % ========== 主流水线 ==========
  \node[stage, fill=mBlue!40] (import)
    {模型导入\\{\footnotesize PyTorch/ONNX}};

  \node[stage, fill=mBlue!40, right=0.8cm of import] (relax)
    {Relax IR\\{\footnotesize 高层次 IR}};

  \node[stage, fill=mBlue!40, right=0.8cm of relax] (te)
    {TE\\{\footnotesize 张量表达式}};

  \node[stage, fill=mGreen!50, right=0.8cm of te] (tuning)
    {MetaSchedule\\{\footnotesize Auto-tuning}};

  \node[stage, fill=mBlue!40, right=0.8cm of tuning] (tir)
    {TIR\\{\footnotesize 底层 IR}};

  \node[stage, fill=mPurple!40, right=0.8cm of tir] (codegen)
    {机器码\\{\footnotesize .so}};

  % ========== 主流程连线 ==========
  \draw[arr] (import) -- (relax)
    node[pos=0.5, above=2pt, ts] {计算图};

  \draw[arr] (relax) -- (te)
    node[pos=0.5, above=2pt, ts] {FuseOps};

  \draw[arr] (te) -- (tuning)
    node[pos=0.5, above=2pt, ts] {schedule};

  \draw[arr] (tuning) -- (tir)
    node[pos=0.5, above=2pt, ts] {JSON DB};

  \draw[arr] (tir) -- (codegen)
    node[pos=0.5, above=2pt, ts] {LLVM};

  % ========== 下方微观细节：每个阶段的关键操作 ==========
  % PyTorch/ONNX 导入
  \node[micro, fill=mBlue!30, below=0.7cm of import] (i1) {torch.onnx};
  \node[micro, fill=mBlue!30, below=0.25cm of i1] (i2) {onnxsim};
  \draw[arr, dashed, thin, mDark!60] (import.south) -- (i1.north);

  % Relax IR 优化
  \node[micro, fill=mBlue!30, below=0.7cm of relax] (r1) {算子融合};
  \node[micro, fill=mBlue!30, below=0.25cm of r1] (r2) {常量折叠};
  \draw[arr, dashed, thin, mDark!60] (relax.south) -- (r1.north);

  % TE 张量表达式
  \node[micro, fill=mBlue!30, below=0.7cm of te] (t1) {循环切分};
  \node[micro, fill=mBlue!30, below=0.25cm of t1] (t2) {矢量化};
  \draw[arr, dashed, thin, mDark!60] (te.south) -- (t1.north);

  % MetaSchedule 调优
  \node[micro, fill=mGreen!40, below=0.7cm of tuning] (m1) {搜索空间};
  \node[micro, fill=mGreen!40, below=0.25cm of m1] (m2) {RPC Runner};
  \node[micro, fill=mGreen!40, below=0.25cm of m2] (m3) {Cost Model};
  \draw[arr, dashed, thin, mDark!60] (tuning.south) -- (m1.north);

  % TIR 底层 IR
  \node[micro, fill=mBlue!30, below=0.7cm of tir] (ti1) {内存重排};
  \node[micro, fill=mBlue!30, below=0.25cm of ti1] (ti2) {循环展开};
  \draw[arr, dashed, thin, mDark!60] (tir.south) -- (ti1.north);

  % Codegen 机器码生成
  \node[micro, fill=mPurple!30, below=0.7cm of codegen] (c1) {AArch64};
  \node[micro, fill=mPurple!30, below=0.25cm of c1] (c2) {cortex-a72};
  \node[micro, fill=mPurple!30, below=0.25cm of c2] (c3) {+neon};
  \draw[arr, dashed, thin, mDark!60] (codegen.south) -- (c1.north);

  % ========== 顶部：端到端流程标注 ==========
  \draw[decorate, decoration={brace, amplitude=4pt, raise=3pt}, thick, draw=mDark]
    (import.north west) -- (codegen.north east)
    node[midway, above=8pt, font=\sffamily\small\bfseries, color=mDark]
    {TVM 端到端编译流水线（baseline $\to$ current）};

  % ========== 张量形状标注 ==========
  \node[ts, below=2.8cm of import, anchor=north] {输入：$1{\times}32{\times}32{\times}32$};
  \node[ts, below=2.8cm of codegen, anchor=north] {输出：$1{\times}3{\times}256{\times}256$};

  % ========== 层次化表达：背景分组 ==========
  \begin{scope}[on background layer]
    % 前端解析与IR优化
    \node[fit=(import)(relax)(te)(i1)(i2)(r1)(r2)(t1)(t2),
      fill=mBlue!12, draw=mDark!40, dashed,
      rounded corners=5pt, inner sep=8pt] {};

    % MetaSchedule 调优
    \node[fit=(tuning)(m1)(m2)(m3),
      fill=mGreen!18, draw=mDark!40, dashed,
      rounded corners=5pt, inner sep=8pt] {};

    % 后端代码生成
    \node[fit=(tir)(codegen)(ti1)(ti2)(c1)(c2)(c3),
      fill=mPurple!12, draw=mDark!40, dashed,
      rounded corners=5pt, inner sep=8pt] {};
  \end{scope}

  % ========== 信息密度拉满：调优配置细节 ==========
  \node[legendbox, anchor=north west] at ([xshift=5pt, yshift=-5pt]current bounding box.north west)
    [align=left, text width=4.8cm] {
    \textbf{MetaSchedule 配置}\\[2pt]
    • 调优策略：baseline-seeded warm-start\\
    • 搜索预算：500 trials（增量）\\
    • Runner：RPC (飞腾派远端)\\
    • Cost Model：XGBoost\\
    • 产物 SHA-256: \texttt{1946b08e...}
  };

  % ========== 图例（右下角）==========
  \node[legendbox, anchor=south east] at ([xshift=-5pt, yshift=5pt]current bounding box.south east)
    [align=left, text width=4.5cm] {
    \textbf{图例}\\[1pt]
    \colorbox{mBlue!40}{\rule{0.16cm}{0.12cm}\hspace{0.05cm}} IR层 \quad
    \colorbox{mGreen!50}{\rule{0.16cm}{0.12cm}\hspace{0.05cm}} 调优层 \quad
    \colorbox{mPurple!40}{\rule{0.16cm}{0.12cm}\hspace{0.05cm}} 代码生成\\[1pt]
    Target: \texttt{\{"kind":"llvm", "mcpu":"cortex-a72",\\
    \quad "mattr":["+neon"], "num-cores":4\}}
  };
\end{tikzpicture}
\end{document}
